\section{Proofs and equivalence checks}
\label{app:b-proofs}

This appendix contains only the proofs needed by the claims of the main text.
The inherited ordinal, abstraction and price-theoretic material that is not
needed is preserved in \texttt{papers/paper\_b/archive/B1\_20260927/}.

\subsection{The population-scoped discrepancy allowance}

\begin{proof}[Proof of \cref{prop:n1-population}]
Let $d$ be a fixed policy.  Write the signed contrast integrand
$h(W)=d(W)\,(U^v(W)-U^0(W))$, with $h(W)\in[-1,1]$ because $U\in[0,1]$ and
$d\in\{0,1\}$.  Let $D$ be the declared joint discrepancy event with
$\Prb(D)\le p$.  Off $D$ the intended and discrepancy-world regime outputs agree,
so the integrand is unchanged.  On $D$, the two regime outcomes can each move by
up to $1$, so for any $W\in D$ the integrand can change by at most $2$; the
contribution of $D$ to the population mean therefore changes by at most
$2\,\Prb(D)\le 2p$.  Combining the unchanged off-$D$ part with the bounded on-$D$
part gives
\[
 \big|\,G_{\text{discrepancy}}(d)-G(d)\,\big|\le 2p .
\]
Hence any valid lower bound $L\le G(d)$ for the intended population gain remains
valid for $G_{\text{discrepancy}}(d)$ after subtracting $2p$, i.e.\
$L-2p\le G_{\text{discrepancy}}(d)$.  The result is a population statement for a
\emph{fixed} policy; selection over a family is handled by simultaneous bounds in
the main text.
\end{proof}

\paragraph{Why the pointwise gate is not available.}
The $2p$ shift is a bound on a population mean, not on a pointwise conditional
mean.  A gate of the form $\mathrm{LCB}(w)-2p>0$ applied per checkpoint would
require a conditional allowance that is not implied by the global event; the
main text therefore forbids it and uses simultaneous population bounds for
family selection.

\subsection{Ridge advantage/action-value equivalence}
\label{app:b-ridge}

Let $B$ be the common ridge fitting operator on paired rows, with identical
design matrices, penalties and rows for the advantage learner and the
action-value learner.  Ridge is affine in the response and the two learners share
the same design matrix, so
\[
 \widehat{\text{adv}} = B\,(y_v-y_0)=B\,y_v-B\,y_0
 = \widehat{q}_v-\widehat{q}_0 ,
\]
up to floating-point error.  On identical designs the two estimators must agree
to within $10^{-12}$; this is an implementation equivalence check, not a
hypothesis and not a contribution.

\subsection{Corrected conditional-expectation statement}
\label{app:b-conditional}

An earlier inherited version of this appendix asserted that, with
cell-measurable policies, the integrand equals the cell-level advantage of the
regime selected in its cell, so that cell-level validity applies
\emph{pointwise}.  That pointwise equality does not hold in general: the
integrand is a random function whose realised value depends on the within-cell
state, whereas the cell-level advantage is a conditional mean.  What holds is
only the \emph{conditional-expectation} identity
\[
 \E\!\left[\,\Delta^{(j)}_{d(Z,M)}(Z,M)\ \middle|\ c(Z,M)\,\right]
 = \Delta^{(j)}_{\text{selected at }c(Z,M)}\big(c(Z,M)\big),
\]
because conditioning on the cell fixes the selected regime.  Cell-level validity
therefore transfers through the conditional expectation, not through a pointwise
equality.  This corrects the inherited assertion.
